\documentclass[11pt,letterpaper]{article}

% ---------- Page layout ----------
\usepackage[letterpaper,margin=0.85in,headheight=14pt]{geometry}
\usepackage{setspace}
\usepackage{microtype}
\usepackage{xcolor}
\usepackage{parskip}
\usepackage{enumitem}
\usepackage{booktabs}
\usepackage{array}
\usepackage{longtable}
\usepackage{fancyhdr}
\usepackage{titlesec}
\usepackage[T1]{fontenc}
\usepackage{lmodern}
\usepackage[hidelinks]{hyperref}

% ---------- Global styling ----------
\definecolor{heading}{gray}{0}
\setstretch{1.08}
\setlength{\parindent}{0pt}
\setlength{\parskip}{5pt}

\setlist[itemize]{
    leftmargin=1.55em,
    itemsep=2pt,
    topsep=3pt,
    parsep=0pt
}

\titleformat{\section}
  {\normalfont\Large\bfseries\color{heading}}
  {\thesection}
  {0.6em}
  {}

\titlespacing*{\section}{0pt}{12pt}{6pt}

\titleformat{\subsection}
  {\normalfont\large\bfseries\color{heading}}
  {\thesubsection}
  {0.6em}
  {}

\titlespacing*{\subsection}{0pt}{9pt}{4pt}

\pagestyle{fancy}
\fancyhf{}
\renewcommand{\headrulewidth}{0pt}
\fancyfoot[C]{\small\thepage}

\newcolumntype{L}[1]{>{\raggedright\arraybackslash}p{#1}}

\begin{document}

% ---------- Title block ----------

\begin{center}
    {\LARGE\bfseries Real-Time Collaborative Code Editor}\par
    \vspace{5pt}
    {\large\itshape Software Engineering Project Proposal}\par
    \vspace{10pt}

    {\normalsize Submitted by: Yuvraj Jasuja(1024030069) and Harwinder(1024030075)}\par
    {\normalsize Department of Computer Engineering}\par
    {\normalsize Thapar Institute of Engineering and Technology}\par
    {\normalsize August 2026}
\end{center}

\vspace{7pt}
\hrule
\vspace{5pt}


% --------------------------------------------------

\section{Higher-order Goal}

The goal of this project is to make collaborative programming more accessible and interactive by providing a shared environment where multiple users can work on the same code in real time.

Programming is often a collaborative activity. Students work together on assignments, developers pair program while solving problems, and teams frequently need to discuss and modify the same codebase. However, basic collaboration often involves repeatedly sharing files, copying code through messaging platforms, or manually merging different versions of the same program.

Our project aims to reduce this friction by creating a real-time coding environment in which users can join a common room and collaborate on the same code simultaneously.


% --------------------------------------------------

\section{Introduction / Problem Statement}

Collaborative coding can become inefficient when multiple people are working on the same programming task.

Consider a common classroom situation where two students are solving a programming problem together. One student writes a part of the solution and sends the file to the other. The second student modifies it and sends another version back. As the number of changes increases, it becomes difficult to keep track of the latest version of the program.

This workflow creates several problems:

\begin{itemize}
    \item Repeated sharing of source code files.
    \item Multiple versions of the same program.
    \item Difficulty identifying the latest version.
    \item Delayed feedback between collaborators.
    \item Version conflicts when different users modify the same code.
    \item Lack of a simple shared environment for collaborative programming.
\end{itemize}

Existing professional collaboration tools can solve many of these problems, but they may be unnecessarily complex for small teams, students, classroom collaboration, or quick coding sessions.

The Real-Time Collaborative Code Editor addresses this problem by allowing multiple users to join a shared coding room and view code changes as they happen.


% --------------------------------------------------

\section{Objectives}

The main objectives of the project are:

\begin{itemize}
    \item Build a web-based platform for real-time collaborative code editing.
    
    \item Allow users to create or join coding rooms using a unique Room ID.
    
    \item Synchronize code changes between multiple connected users in real time.
    
    \item Display active users currently connected to a coding room.
    
    \item Provide a clean and interactive code editing environment.
    
    \item Ensure that users joining an existing room can receive the current state of the code.
    
    \item Extend the platform in future phases to support multiple programming languages and secure code execution.
\end{itemize}




% --------------------------------------------------

\section{Methodology / Solution Approach}

\subsection{System Architecture}

The system follows a client-server architecture with real-time communication between users and the backend server.

The main components are:

\begin{itemize}
    \item User Interface built using React.
    
    \item Code Editor implemented using CodeMirror.
    
    \item Real-time communication handled using Socket.IO.
    
    \item Backend server developed using Node.js and Express.
    
    \item Room-based communication to separate different collaborative coding sessions.
\end{itemize}

The overall architecture can be represented as:

\begin{center}
\textbf{
User A $\rightarrow$
React + CodeMirror $\rightarrow$
Socket.IO $\rightarrow$
Node.js Server
}
\end{center}

The server then broadcasts relevant changes to other users connected to the same room:

\begin{center}
\textbf{
Node.js Server $\rightarrow$
Socket.IO $\rightarrow$
User B, User C, ...
}
\end{center}


\subsection{Real-Time Collaboration Flow}

The main workflow of the system is as follows:

\begin{enumerate}
    \item A user enters a username and creates a new coding room or enters an existing Room ID.
    
    \item The user establishes a Socket.IO connection with the backend server.
    
    \item The server associates the user with the requested coding room.
    
    \item Other users can join the same room using the same Room ID.
    
    \item When a user modifies the code, the editor detects the change.
    
    \item The updated code is emitted to the backend using Socket.IO.
    
    \item The server broadcasts the update to other users in the same room.
    
    \item Connected users receive the update and their editor reflects the new code.
    
    \item When a user disconnects, the server notifies the remaining users.
\end{enumerate}


\subsection{Room-Based Communication}

Each collaborative session is organized as a separate room.

This prevents code changes from one group of users from being sent to unrelated users.

For example:

\begin{itemize}
    \item Room A may contain User 1 and User 2.
    \item Room B may contain User 3 and User 4.
\end{itemize}

Code changes generated inside Room A are only broadcast to users connected to Room A.

This makes the system more organized and allows multiple independent coding sessions to run simultaneously.


\subsection{Future Multi-Language Execution}

The current version of the project focuses primarily on real-time collaborative editing.

A future extension of the platform will allow users to select programming languages such as:

\begin{itemize}
    \item C++
    \item Java
    \item Python
    \item JavaScript
    \item C
\end{itemize}

The user will be able to write code, provide input, execute the program, and view output or compilation errors.

Because executing arbitrary user code directly on the application server can create security risks, code execution will be handled through a secure sandbox or external code-execution service.

The future execution pipeline will follow:

\begin{center}
\textbf{
Code Editor $\rightarrow$
Execution API / Secure Sandbox $\rightarrow$
Compiler / Runtime $\rightarrow$
Output
}
\end{center}


% --------------------------------------------------

\section{Tech Stack}

The project uses modern web technologies suitable for real-time applications.

\begin{itemize}
    \item \textbf{React:} Used to build the frontend user interface and manage application components.
    
    \item \textbf{CodeMirror:} Used to provide an interactive code editing environment.
    
    \item \textbf{Socket.IO:} Used for real-time, bidirectional communication between clients and the server.
    
    \item \textbf{Node.js:} Used to run the backend application.
    
    \item \textbf{Express.js:} Used to create and manage backend routes and server functionality.
    
    \item \textbf{UUID:} Used for generating unique Room IDs.
    
    \item \textbf{Render:} Can be used for deploying the application as a web service.
\end{itemize}


% --------------------------------------------------

\newpage

\section{Evaluation Criteria}

The success of the project will be evaluated using practical and measurable criteria.

\begin{itemize}

    \item \textbf{Synchronization Accuracy:} Whether code changes made by one user are correctly reflected for other users in the same room.

    \item \textbf{Real-Time Response:} The delay between a code change made by one user and the update received by another user.

    \item \textbf{Room Isolation:} Verification that users in different rooms do not receive each other's code updates.

    \item \textbf{Connection Handling:} Correct handling of users joining and leaving a room.

    \item \textbf{Concurrent Users:} The ability of the application to support multiple users within a shared coding room.

    \item \textbf{Usability:} Whether users can easily create a room, join a room, and begin collaborating.
    
\end{itemize}

A basic experiment can be performed by opening the same Room ID on multiple browsers or devices and checking whether changes made by one user are reflected correctly for the others.


% --------------------------------------------------

\section{Scalability}

The architecture is designed around independent real-time rooms.

As the number of users increases, the application can support multiple collaborative sessions by maintaining separate room connections.

Potential scalability improvements include:

\begin{itemize}
    \item Deploying multiple backend instances.
    
    \item Using a distributed Socket.IO adapter for communication between server instances.
    
    \item Adding persistent storage for projects and code files.
    
    \item Using a database to manage users and saved projects.
    
    \item Separating the code execution service from the main collaboration server.
\end{itemize}

This approach would allow the project to evolve from a small collaborative editor into a larger collaborative coding platform.


% --------------------------------------------------

\section{Resources / Tooling}

The project is designed using widely available and mostly open-source technologies.

\begin{itemize}
    \item React and JavaScript for frontend development.
    
    \item CodeMirror for the coding interface.
    
    \item Node.js and Express for backend development.
    
    \item Socket.IO for real-time communication.
    
    \item GitHub for version control and project management.
    
    \item Render or another cloud platform for deployment.
    
    \item A secure execution API or sandbox for future multi-language code execution.
\end{itemize}

The use of open-source technologies allows the project to be developed without significant infrastructure costs.


% --------------------------------------------------

\newpage

\section{Project Scope and Deliverables}

\subsection{Initial Deliverable}

The first version of the project will include:

\begin{itemize}
    \item Home page for entering a username and Room ID.
    
    \item Ability to create a new unique coding room.
    
    \item Ability to join an existing coding room.
    
    \item Real-time collaborative code editing.
    
    \item Display of connected users.
    
    \item Room-based communication using Socket.IO.
    
    \item User connection and disconnection handling.
\end{itemize}


\subsection{Subsequent Deliverables}

Future versions of the project may include:

\begin{itemize}
    \item Multiple programming language support.
    
    \item Syntax highlighting based on selected language.
    
    \item Run code functionality.
    
    \item Custom input and output terminal.
    
    \item Compilation and runtime error display.
    
    \item Execution time information.
    
    \item User authentication.
    
    \item Saving and loading projects.
    
    \item Multi-file project support.
    
    \item Live cursor tracking.
    
    \item Integrated chat for collaborators.
\end{itemize}


% --------------------------------------------------

\section{Risks and Mitigations}

\begin{itemize}

    \item \textbf{Synchronization conflicts:}
    Multiple users may modify the code at nearly the same time. The initial version relies on event-based synchronization, while more advanced versions can introduce operational transformation or CRDT-based collaboration.

    \item \textbf{Network latency:}
    Slow or unstable internet connections may delay updates. Socket.IO provides reconnection support to improve reliability.

    \item \textbf{User disconnection:}
    A user may unexpectedly lose connection. The application handles disconnect events and updates the active user list.

    \item \textbf{Unauthorized room access:}
    Future versions can introduce authentication, private rooms, and permission-based access.

    \item \textbf{Security of code execution:}
    Running arbitrary user code directly on the main server can be dangerous. Future execution features will use a secure sandbox or isolated execution service.

\end{itemize}


% --------------------------------------------------

\section{Timeline}

\begin{center}
\small

\begin{tabular}{@{}L{0.18\linewidth} L{0.76\linewidth}@{}}

\toprule

\textbf{Phase} & \textbf{Focus} \\

\midrule

Weeks 1--2 &
Requirement analysis, project setup, React frontend, and room creation interface. \\

Weeks 3--4 &
CodeMirror integration and basic code editing functionality. \\

Weeks 5--6 &
Socket.IO integration and real-time synchronization between users. \\

Weeks 7--8 &
Room management, connected users, and join/disconnect handling. \\

Weeks 9--10 &
Testing with multiple users, UI improvements, and deployment. \\

Weeks 11--12 &
Future feature prototype including language selection and secure code execution integration. \\

\bottomrule

\end{tabular}

\end{center}


% --------------------------------------------------

\section{Summary}

The Real-Time Collaborative Code Editor is designed to simplify collaborative programming by providing a shared environment where multiple users can work on the same code simultaneously.

The platform combines React and CodeMirror for the frontend coding experience with Node.js, Express, and Socket.IO for real-time communication. Users can create or join coding rooms, collaborate with other users, and see code changes as they happen.

The current project focuses on solving the core challenge of real-time code synchronization. Future development will extend the platform with multi-language support, secure code execution, project storage, authentication, and additional collaborative features.

Ultimately, the project aims to evolve from a basic real-time editor into a practical collaborative coding environment for students, developers, and teams.

\end{document}
